\documentclass[11pt,a4paper]{article}

\usepackage[brazil]{babel}
\usepackage[utf8]{inputenc}
\usepackage[T1]{fontenc}
\usepackage{lmodern}
\usepackage{geometry}
\usepackage{setspace}
\usepackage{parskip}
\usepackage{titlesec}
\usepackage{xcolor}
\usepackage{hyperref}

\geometry{margin=2.5cm}
\setstretch{1.2}

\definecolor{primary}{RGB}{41,128,185}
\hypersetup{
    colorlinks=true,
    linkcolor=primary,
    urlcolor=primary
}

\titleformat{\section}{\Large\bfseries\color{primary}}{\thesection}{1em}{}
\titleformat{\subsection}{\large\bfseries}{\thesubsection}{1em}{}

\title{\textbf{Composição da Chapa 01 \\ GULECD/UFPR}}
\author{Chapa 01}
\date{\today}

\begin{document}

\maketitle

\section{Introdução}

Este documento apresenta a composição completa da chapa 01 candidata à coordenação do GULECD/UFPR, incluindo os cargos propostos e uma breve descrição do perfil de cada integrante.

\section{Coordenação Representativa}

\subsection{Coordenador Representativo}

\textbf{Nome:} Marcos Oliveira de Carvalho

\textbf{Curso/Vínculo:} Estatística e Ciência de Dados

\textbf{Mini-biografia:}  
Marcos de Carvalho é farmacêutico-bioquímico (UFSM), com mestrado em Biologia Celular e Molecular (UFRGS), doutorado em Genética e Biologia Molecular (UFRGS) e formação em empreendedorismo (Babson College). Possui experiência em bioinformática e biologia computacional, com atuação em áreas como genômica, metagenômica e oncogenômica. É fundador de duas startups de biotecnologia, com trajetória voltada à inovação e desenvolvimento de soluções baseadas em dados integradas a processos laboratoriais. Atualmente, atua como pesquisador de pós-doutorado na área de medicina translacional no CHC-UFPR, com interesse em integração de dados clínicos e moleculares, ao mesmo tempo em que cursa graduação em Estatística e Ciência de Dados (UFPR). Tem grande interesse em software livre/código aberto, ciência/tecnologia e no desenvolvimento de infraestrutura técnica para pesquisa e formação acadêmica. É usuário de sistemas Linux desde 1999, com experiência desde computadores pessoais e servidores até clusters de computação de alta performance.

\vspace{1em}

\subsection{Vice-Coordenador Representativo}

\textbf{Nome:} [Inserir Nome Completo]

\textbf{Curso/Vínculo:} [Inserir curso, programa ou vínculo institucional]

\textbf{Mini-biografia:}  
[Inserir uma breve descrição do membro, incluindo formação, experiência relevante, participação no grupo e interesses acadêmicos ou profissionais.]

\section{Núcleo de Coordenação}

\subsection{Membro do Núcleo de Coordenação 1}

\textbf{Nome:} [Inserir Nome Completo]

\textbf{Cargo Proposto:} Coordenador Associado

\textbf{Curso/Vínculo:} [Inserir curso, programa ou vínculo institucional]

\textbf{Mini-biografia:}  
[Inserir uma breve descrição do membro, incluindo formação, experiência relevante, participação no grupo e interesses acadêmicos ou profissionais.]

\vspace{1em}

\subsection{Membro do Núcleo de Coordenação 2}

\textbf{Nome:} [Inserir Nome Completo]

\textbf{Cargo Proposto:} Coordenador Associado

\textbf{Curso/Vínculo:} [Inserir curso, programa ou vínculo institucional]

\textbf{Mini-biografia:}  
[Inserir uma breve descrição do membro, incluindo formação, experiência relevante, participação no grupo e interesses acadêmicos ou profissionais.]

\vspace{1em}

\subsection{Membro do Núcleo de Coordenação 3}

\textbf{Nome:} [Inserir Nome Completo]

\textbf{Cargo Proposto:} Coordenador Associado

\textbf{Curso/Vínculo:} [Inserir curso, programa ou vínculo institucional]

\textbf{Mini-biografia:}  
[Inserir uma breve descrição do membro, incluindo formação, experiência relevante, participação no grupo e interesses acadêmicos ou profissionais.]

\vspace{1em}

\subsection{Membro do Núcleo de Coordenação 4}

\textbf{Nome:} [Inserir Nome Completo]

\textbf{Cargo Proposto:} Coordenador Associado

\textbf{Curso/Vínculo:} [Inserir curso, programa ou vínculo institucional]

\textbf{Mini-biografia:}  
[Inserir uma breve descrição do membro, incluindo formação, experiência relevante, participação no grupo e interesses acadêmicos ou profissionais.]

\vspace{1em}

\subsection{Membro do Núcleo de Coordenação 5}

\textbf{Nome:} [Inserir Nome Completo]

\textbf{Cargo Proposto:} Coordenador Associado

\textbf{Curso/Vínculo:} [Inserir curso, programa ou vínculo institucional]

\textbf{Mini-biografia:}  
[Inserir uma breve descrição do membro, incluindo formação, experiência relevante, participação no grupo e interesses acadêmicos ou profissionais.]

\vspace{1em}

\subsection{Membro do Núcleo de Coordenação 6}

\textbf{Nome:} [Inserir Nome Completo]

\textbf{Cargo Proposto:} Coordenador Associado

\textbf{Curso/Vínculo:} [Inserir curso, programa ou vínculo institucional]

\textbf{Mini-biografia:}  
[Inserir uma breve descrição do membro, incluindo formação, experiência relevante, participação no grupo e interesses acadêmicos ou profissionais.]

\vspace{1em}

\subsection{Membro do Núcleo de Coordenação 7}

\textbf{Nome:} [Inserir Nome Completo]

\textbf{Cargo Proposto:} Coordenador Associado

\textbf{Curso/Vínculo:} [Inserir curso, programa ou vínculo institucional]

\textbf{Mini-biografia:}  
[Inserir uma breve descrição do membro, incluindo formação, experiência relevante, participação no grupo e interesses acadêmicos ou profissionais.]

\vspace{1em}

\subsection{Membro do Núcleo de Coordenação 8}

\textbf{Nome:} [Inserir Nome Completo]

\textbf{Cargo Proposto:} Coordenador Associado

\textbf{Curso/Vínculo:} [Inserir curso, programa ou vínculo institucional]

\textbf{Mini-biografia:}  
[Inserir uma breve descrição do membro, incluindo formação, experiência relevante, participação no grupo e interesses acadêmicos ou profissionais.]

\section{Considerações Finais}

A composição desta chapa foi estruturada com o objetivo de garantir diversidade de competências, equilíbrio entre experiência e potencial de desenvolvimento, e alinhamento com os princípios do GULECD/UFPR.

Todos os membros comprometem-se com a atuação colaborativa, ética e transparente na condução das atividades do grupo.

\end{document}